% SOURCE_AUTHORITY: sga5_fr_workpass.tex, lines 14575--14603 (LF-normalized unit).
% SOURCE_SHA256: 791F4EFFC5E02832D5D77ED03518C8156D6F07E4C8238B03545DB93D883FBB28
Naturalmente, se define del mismo modo una clase de correspondencia de Frobenius iterada $(\fr_X^\nu,(\Fr^*_{F/X})^\nu)$ para todo entero $\nu\geq0$; el isomorfismo
\[
(\Fr^*_{F/X})^\nu:(\fr_X^\nu)^*(F)\longrightarrow F
\]
es igual a
\[
(\Fr^*_{F/X})^{\nu-1}\circ (\fr_X^{\nu-1})^*(\Fr^*_{F/X})
\]
o también a
\[
\Fr^*_{F/X}\circ\fr_X^*((\Fr^*_{F/X})^{\nu-1})
\]
para $\nu\geq1$.

Examinemos el caso particular en que $F$ es representable por un preesquema $V$ étale sobre $X$; entonces $(\fr_X)^*(F)=V^{(p/X)}$, $F(V)=\Hom_X(V,V)$,
\[
(\fr_X)^*(F)(V)=F(V^{(p/X)})=\Hom_X(V^{(p/X)},V),
\]
y $F(\Fr_{V/X})$ es la aplicación de $\Hom_X(V^{(p/X)},V)$ en $\Hom_X(V,V)$ que transforma un morfismo $\varphi$ en $\varphi\circ\Fr_{V/X}$; ahora bien, $\Fr^*_{F/X}$ se identifica con un morfismo de $V^{(p/X)}$ en $V$ que es la imagen de $\id_V$ por la aplicación $F(\Fr_{V/X})^{-1}$ inversa de la precedente. Esto muestra que, para $F=V$ étale sobre $X$, se tiene
\[
\Fr^*_{V/X}=(\Fr_{V/X})^{-1}.
\]

Obsérvese que, si el haz $F$ es constante, es claro que
\[
(\fr_X)^*(F)=F,
\qquad
\Fr^*_{F/X}=\id_F.
\]
